% P11 paper module: R26 concrete near-null survival of relative resolvents.

\subsection{Near-null survival of the concrete relative resolvents}
\label{subsec:o3y-nearnull-resolvent}

R24 and R25 reduce the nested inverse-root problem to resolvents of the concrete
relative metrics.  The absolute odd terminal forms nevertheless diverge on every fixed
nonzero smooth odd vector.  The next result shows that this absolute divergence does
\emph{not} drive the relative resolvents to zero: the R17 near-null directions leave a
nonvanishing inverse-metric trace on every fixed higher-jet direction.

Fix $0<R<S<T_0<U$.  Write
\[
 B_R:=G_{R,T_0}^-,
 \qquad
 C_R(U):=G_{R,U}^-,
 \qquad
 A_R(U):=B_R^{-1/2}C_R(U)B_R^{-1/2},
\tag{O3Y.1}\label{eq:o3y-relative-R}
\]
and analogously $B_S,C_S(U),A_S(U)$.  The normalized base transition is
\[
 W=B_S^{1/2}J_{R,S}B_R^{-1/2}.
\]

\begin{theorem}[Higher-jet fixed-vector resolvent survival]
\label{thm:o3y-resolvent-survival}
Let
\[
 0\ne f\in C_c^\infty((-R,R))_{\rm odd},
 \qquad
 \beta_R^{(0)}(f)=0.
\]
Put
\[
 x_f:=B_R^{1/2}f,
 \qquad
 b_f:=\|x_f\|_{X,R}^2
 =q_{T_0}^X(J_{R,T_0}f),
 \qquad
 \gamma_f:=\mathfrak c_{\Gamma,R}[f]>0.
\]
Then for every fixed $t\ge0$,
\[
 \boxed{
 \liminf_{U\to\infty}
 \left\langle(A_R(U)+tI)^{-1}x_f,x_f\right\rangle_{X,R}
 \ge
 \frac{b_f^2}{\gamma_f+t b_f}>0.
 }
\tag{O3Y.2}\label{eq:o3y-R-resolvent-lower}
\]
At the canonically nested source level,
\[
 \boxed{
 \liminf_{U\to\infty}
 \left\langle(A_S(U)+tI)^{-1}Wx_f,Wx_f\right\rangle_{X,S}
 \ge
 \frac{b_f^2}{\gamma_f+t b_f}>0.
 }
\tag{O3Y.3}\label{eq:o3y-S-resolvent-lower}
\]
In particular neither relative resolvent family converges strongly to zero on these
fixed vectors.
\end{theorem}

\begin{proof}
By Corollary~\ref{cor:first-jet}, $f$ has a finite first nonzero integral-jet order
$m>0$.  Choose a fixed smooth odd $f_0$ with first jet order zero.  Proposition
\ref{prop:o3n-nearnull} defines
\[
 z_U
 =f-\frac{\ell_U(f)}{\ell_U(f_0)}f_0,
 \qquad
 \frac{\ell_U(f)}{\ell_U(f_0)}=O(U^{-m}),
\]
so $z_U\to f$ in the fixed smooth two-dimensional source span and hence in the fixed
$X$-graph norm.  Corollary~\ref{cor:o3p-gamma-limit} gives
\[
 q_U^X(J_{R,U}z_U)\longrightarrow\gamma_f.
\tag{O3Y.4}\label{eq:o3y-nearnull-energy}
\]

Set
\[
 y_U:=B_R^{1/2}z_U.
\]
Then $y_U\to x_f$ and, by the definition of the relative metric,
\[
 \langle A_R(U)y_U,y_U\rangle_{X,R}
 =q_U^X(J_{R,U}z_U)\longrightarrow\gamma_f.
\tag{O3Y.5}\label{eq:o3y-relative-nearnull}
\]
Also
\[
 \langle x_f,y_U\rangle_{X,R}\to b_f,
 \qquad
 \|y_U\|_{X,R}^2\to b_f.
\tag{O3Y.6}\label{eq:o3y-base-limits}
\]

For every positive invertible $A$ and $t\ge0$,
\[
 \langle(A+tI)^{-1}x,x\rangle
 =\sup_{y\ne0}
 \frac{|\langle x,y\rangle|^2}
 {\langle Ay,y\rangle+t\|y\|^2}.
\tag{O3Y.7}\label{eq:o3y-variational}
\]
Apply this with $A=A_R(U)$, $x=x_f$, and the admissible test vector $y=y_U$.
Equations~\eqref{eq:o3y-relative-nearnull}--\eqref{eq:o3y-base-limits} give
\eqref{eq:o3y-R-resolvent-lower}.

For level $S$, R17 source compatibility gives
$z_U^S=J_{R,S}z_U$.  Hence its standardized representative is
\[
 B_S^{1/2}z_U^S
 =B_S^{1/2}J_{R,S}B_R^{-1/2}y_U
 =Wy_U,
\]
while the fixed representative is $Wx_f$.  Since $W$ is an isometry, the numerator and
base norms have the same limits as at level $R$.  The terminal cocycle gives
\[
 \langle A_S(U)Wy_U,Wy_U\rangle_{X,S}
 =q_U^X(J_{S,U}J_{R,S}z_U)
 =q_U^X(J_{R,U}z_U)
 \to\gamma_f.
\]
A second application of~\eqref{eq:o3y-variational} proves
\eqref{eq:o3y-S-resolvent-lower}.
\end{proof}

\begin{corollary}[Concrete inverse-root survival]
\label{cor:o3y-inverse-root-survival}
Under the hypotheses of Theorem~\ref{thm:o3y-resolvent-survival},
\[
 \boxed{
 \liminf_{U\to\infty}
 \|A_R(U)^{-1/2}x_f\|_{X,R}^2
 \ge\frac{b_f^2}{\gamma_f}>0,
 }
\tag{O3Y.8}\label{eq:o3y-R-inverse-lower}
\]
and
\[
 \boxed{
 \liminf_{U\to\infty}
 \|A_S(U)^{-1/2}Wx_f\|_{X,S}^2
 \ge\frac{b_f^2}{\gamma_f}>0.
 }
\tag{O3Y.9}\label{eq:o3y-S-inverse-lower}
\]
Thus $A_R(U)^{-1/2}$ and $A_S(U)^{-1/2}$ do not converge strongly to zero on the
higher-jet fixed vectors above.
\end{corollary}

\begin{proof}
Take $t=0$ in Theorem~\ref{thm:o3y-resolvent-survival} and use
\[
 \langle A^{-1}x,x\rangle=\|A^{-1/2}x\|^2.
\]
\end{proof}

\begin{remark}[Concrete nonuniformity of the odd terminal metric]
\label{rem:o3y-nonuniformity}
For the same fixed $f$ with first jet order $m>0$, the sharp odd asymptotic gives
\[
 q_U^X(J_{R,U}f)
 \asymp_f\frac{e^U}{U^{2m+2}}
 \longrightarrow\infty.
\]
Equivalently, the relative quadratic form diverges on the fixed standardized vector
$x_f$.  Yet the resolvent and inverse-root quantities in
\eqref{eq:o3y-R-resolvent-lower} and~\eqref{eq:o3y-R-inverse-lower} stay bounded away
from zero.

There is no contradiction.  The moving vector
\[
 z_U=f+O(U^{-m})f_0
\]
converges to $f$ in the fixed source norm while cancelling the dominant constant-mode
boundary functional exactly.  Thus the future forms are strongly nonuniform in
orientation.  In particular, pointwise divergence of the future quadratic form on the
fixed smooth odd core cannot be promoted to strong extinction of its inverse functional
calculus.
\end{remark}

\begin{corollary}[Strong-resolvent extinction is excluded for the actual odd P11 family]
\label{cor:o3y-no-extinction}
For every fixed $t\ge0$, the concrete families
\[
 (A_R(U)+tI)^{-1}
 \quad\text{and}\quad
 (A_S(U)+tI)^{-1}
\]
do not converge strongly to zero on the odd source spaces.
\end{corollary}

\begin{proof}
Choose any nonzero smooth odd $f$ with $\beta_R^{(0)}(f)=0$ and apply
Theorem~\ref{thm:o3y-resolvent-survival}.  A strong zero limit would force the displayed
positive matrix element to tend to zero, contradicting
\eqref{eq:o3y-R-resolvent-lower} or~\eqref{eq:o3y-S-resolvent-lower}.
\end{proof}

\begin{remark}[R26 firewall: survival of both terms does not decide their difference]
\label{rem:o3y-firewall}
R25 introduced the positive resolvent compression defect
\[
 \mathscr E_U(t)
 =W^*(A_S(U)+tI)^{-1}W-(A_R(U)+tI)^{-1}\ge0.
\]
Theorem~\ref{thm:o3y-resolvent-survival} shows that on every fixed smooth higher-jet
vector both resolvent terms have positive liminf.  Therefore R25-F cannot be settled by
a crude claim that the two resolvents separately vanish on the smooth odd core.

The theorem does not determine their difference: it gives neither
\[
 \langle\mathscr E_U(t)x_f,x_f\rangle\to0
\]
nor a positive limsup.  The remaining R25 question is genuinely relative: how much
additional inverse-metric relaxation is gained by allowing the level-$S$ complement of
$\Ran W$?  This result remains on the metric/resolvent layer and gives no conclusion
about $\Gamma_U$, the R22 angle defect, strong terminal transport, or a global
Object~X.
\end{remark}
